\chapter{Estrategia empírica}
\label{chap:estrategia}

% ==========================================================================
% CAMBIO RESPECTO A LA VERSIÓN 2024
% 2024: unidad código postal-mes; instrumento T_i = "infracciones por encima
%       de la mediana hasta un año antes"; DiD + IV; forma reducida
%       ln(acc+1) ~ T_i + Post_t + T_i·Post_t.
% Ahora: unidad colonia-mes; tratamiento = exposición espacial a la
%       infraestructura de fiscalización; estudio de evento mes a mes; sin IV
%       (las multas entran como mecanismo, no como primera etapa).
% ==========================================================================

\section{Evento y unidad de análisis}

El cambio de política que explota este trabajo es la \textbf{renovación de la
Policía de Tránsito de la Ciudad de México de abril de 2022}: la creación de un
cuerpo especializado, con procesos digitalizados y agentes facultados para
infraccionar. A diferencia de la versión previa de esta investigación, que
tomaba la reforma del Reglamento de Tránsito como un cambio homogéneo en toda la
ciudad, aquí el interés está en su \emph{implementación diferencial en el
espacio}: la hipótesis es que una reforma de fiscalización se aplica con más
intensidad donde ya existe infraestructura operativa.

La unidad de observación es la \textbf{colonia-mes}. Se utiliza la división en
unidades territoriales del Instituto Electoral de la Ciudad de México
(1{,}814 unidades). La unidad se define administrativamente y es independiente
de si la colonia resulta tratada o no. Los \emph{outcomes} se agregan a este
mismo nivel.

\section{Definición del tratamiento}

Para cada colonia $i$ se construye una medida de \textbf{exposición a la
infraestructura de fiscalización} (radares fijos y móviles, y equipos
asociados) a partir del inventario del año anterior al evento (2021). Se
consideran tres definiciones, que se reportan como análisis de sensibilidad al
\emph{buffer}:

\begin{equation}
\text{tratada}_i(b) = \mathbbm{1}\{\text{existe al menos un equipo a menos de } b \text{ metros de la colonia } i\},
\qquad b \in \{250, 500, 1000\}.
\end{equation}

Como robustez se emplean también medidas continuas de exposición: el número de
equipos dentro del \emph{buffer} y la proporción del área de la colonia dentro
de $b$ metros de algún equipo. El número de colonias tratadas va de
aproximadamente 215 ($b=250$ m) a 590 ($b=1000$ m).

\section{Especificación}

Para el evento del mes $E$ se estima un estudio de evento mes a mes:

\begin{equation}
\label{eq:event_study}
y_{it} \;=\; \alpha_i \;+\; \lambda_t \;+\;
\sum_{k \neq -1} \beta_k \,\bigl(\mathbbm{1}\{\tau_{it} = k\} \times \text{tratada}_i\bigr)
\;+\; \varepsilon_{it},
\end{equation}

donde $\tau_{it}$ es el número de meses entre el mes $t$ y el evento $E$ (de modo
que $\tau = 0$ en abril de 2022), $k$ recorre esos valores omitiendo $k=-1$
(marzo de 2022) como periodo de referencia, $\alpha_i$ son efectos fijos de
colonia, $\lambda_t$ efectos fijos de mes-año, y $y_{it}$ el \emph{outcome} de
interés. Los errores estándar se agrupan a nivel colonia.

La ventana del análisis es de $-10$ a $+10$ meses. El límite inferior evita el
periodo con problemas de cobertura en la base de infracciones (abril a junio de
2021) y mantiene el análisis fuera de la fase aguda de la pandemia. El límite
superior es el máximo que permite la fuente de infracciones para el evento de
2022.

Además del estudio de evento se reporta un efecto promedio posterior mediante
una diferencia en diferencias agrupada:

\begin{equation}
\label{eq:did}
y_{it} \;=\; \alpha_i \;+\; \lambda_t \;+\; \delta \,\bigl(\text{tratada}_i \times \text{Post}_t\bigr) \;+\; \varepsilon_{it},
\end{equation}

con $\text{Post}_t = \mathbbm{1}\{\tau_{it} \geq 0\}$.

\section{Outcomes}

\begin{itemize}
  \item \textbf{Principal.} Víctimas de hechos de tránsito registradas en
        carpetas de investigación de la Fiscalía General de Justicia: lesiones y
        homicidios culposos por tránsito vehicular. Se analizan por separado el
        total de víctimas y las víctimas fallecidas. Esta fuente cubre hasta
        julio de 2024, lo que permite observar hasta $+10$ meses después del
        evento.
  \item \textbf{Secundario.} Incidentes viales reportados al C5. Se reporta con
        la advertencia de que, para este \emph{outcome}, las tendencias previas
        no son planas (Sección~\ref{sec:pretendencias}).
  \item \textbf{Mecanismo.} Infracciones de tránsito, totales y de exceso de
        velocidad (artículo 9). No constituyen una primera etapa ni se utiliza
        un estimador de variables instrumentales: sirven únicamente para
        documentar si la reforma efectivamente aumentó la actividad de
        fiscalización de forma diferencial cerca de la infraestructura. Un
        eventual efecto disuasivo sobre las víctimas puede operar sin pasar por
        la emisión de multas.
\end{itemize}

\section{Supuesto de identificación}

La identificación descansa en el supuesto de \textbf{tendencias paralelas}: en
ausencia de la reforma, la evolución del \emph{outcome} en las colonias
expuestas a la infraestructura y en las no expuestas habría sido la misma. Este
supuesto se evalúa con la prueba conjunta de que los coeficientes previos al
evento ($\beta_k$ con $k < -1$) son cero (prueba de Wald, 9 restricciones,
errores agrupados por colonia), y visualmente con el propio estudio de evento.

\section{Amenazas a la identificación}

\begin{enumerate}
  \item \textbf{Ubicación no aleatoria de la infraestructura.} Los equipos de
        fiscalización se colocan en vías primarias, que tienen una dinámica de
        tráfico propia. Si esa dinámica evoluciona distinto que en el resto de
        la ciudad, se viola el supuesto de tendencias paralelas. Esto es
        precisamente lo que ocurre con el \emph{outcome} de incidentes C5.
  \item \textbf{Efectos de reporte.} Una mayor presencia policial puede
        aumentar la probabilidad de que una lesión de tránsito se registre
        formalmente como carpeta de investigación, sin que haya un cambio real
        en la siniestralidad. Las víctimas fallecidas son mucho menos sensibles
        a este canal.
  \item \textbf{Otros cambios simultáneos.} Se consideran otros eventos de
        política vial alrededor de abril de 2022 (Sección~\ref{sec:contexto}).
\end{enumerate}
